\documentclass[11pt]{article}
\usepackage[T1]{fontenc}
\usepackage[utf8]{inputenc}
\usepackage{lmodern,microtype,geometry,amsmath,amssymb,amsthm,mathtools}
\usepackage[colorlinks=true,linkcolor=blue,urlcolor=blue]{hyperref}
\hypersetup{pdftitle={Few online mistakes imply a small Boolean circuit for explicit circuit-induction rules},pdfauthor={ChatGPT-6 Astra; Kaarel Hanni; Dmitry Vaintrob}}
\geometry{margin=1in}
\newtheorem{theorem}{Theorem}
\newtheorem{lemma}[theorem]{Lemma}
\newtheorem{proposition}[theorem]{Proposition}
\newtheorem{corollary}[theorem]{Corollary}
\newtheorem{remark}[theorem]{Remark}
\newcommand{\E}{\mathbb E}
\newcommand{\Prb}{\mathbb P}
\newcommand{\cC}{\mathcal C}
\newcommand{\cX}{\mathcal X}
\newcommand{\cH}{\mathcal H}
\newcommand{\1}{\mathbf 1}
\title{Few online mistakes imply a small Boolean circuit\\for explicit circuit-induction rules}
\author{ChatGPT-6 Astra \and Kaarel H\"anni \and Dmitry Vaintrob\thanks{Authors are listed in alphabetical order.}}
\date{21 October 2026}
\begin{document}
\maketitle
\begin{abstract}
A small circuit gives an online mistake bound under an appropriate circuit
prior. We prove converses for three concrete learning rules. A learner that
always selects a smallest consistent circuit makes at least
$C(f)/(3d+3)$ mistakes before learning a Boolean function on its full
$d$-dimensional domain. For Bayesian posterior sampling with a geometric
gate-count prior and uniform wiring, expected total mistakes $M$ imply
$C(f)=O(dM\log(d+dM))$. A prefix circuit language with explicit exception
constructors improves this to $O(dM+1)$ while retaining the usual
$O(C(f)\log(d+C(f))+1)$ upper bound. The Bayesian proofs control logarithmic
loss by sampling mistakes using injective, cheap repairs of circuit
descriptions. They apply to any order revealing the whole domain and to
infinite iid sampling on any fixed support. Coverage, the specified learning
rule, and posterior sampling are substantive hypotheses. The conclusions
do not hold for arbitrary learners or arbitrary priors, and the Bayesian
argument gives no converse for majority prediction. These are mathematical
results, not Lean-verified formalizations.
\end{abstract}

\section{Setup and what a converse can say}

Fix $d\ge1$ and $\cX=\{0,1\}^d$. A deterministic target is
$f:\cX\to\{0,1\}$. A circuit consists of a topologically ordered list of
two-input NAND gates. Its available initial wires are the $d$ input bits
and constants $0,1$. Each gate chooses an ordered pair of earlier wires,
with repetitions allowed. The output is the last gate; the empty list
outputs the constant zero. Dead gates are permitted. Its size $|c|$ is its
number of gates. Let
\[
 C_A(f)=\min\{|c|:c(x)=f(x)\text{ for every }x\in A\},
 \qquad C(f)=C_{\cX}(f).
\]
All these minima exist. Our conclusions concern total NAND gate count,
with no fixed-depth restriction.

At round $t$, an input $x_t$ is selected, the learner predicts, and the
correct label $f(x_t)$ is revealed. Repetitions have the same label.
The input may depend on past predictions and labels, but it must be
selected before the fresh random prediction on that round. Write
$A_t=\{x_1,\ldots,x_t\}$. A complete reveal of $A$ means that precisely
points from $A$ have been queried and every point of $A$ has appeared at
least once. Inputs outside $A$ need not be learned.

Coverage is unavoidable: labels on a small set of easy inputs cannot
determine the target's complexity elsewhere. Likewise, an arbitrary
learner may hard-code a highly complex target and make no mistakes. The
theorems below concern specified circuit-induction procedures and their
complete observations, not a representation theorem for every successful
predictor.

\section{A one-point circuit repair}

Put $R_d=3d+3$.

\begin{lemma}[Injective overwrite]\label{lem:repair}
For each $(x,y)\in\cX\times\{0,1\}$ there is a map
$P_{x,y}:\cC\to\cC$ that appends exactly $r_x\le R_d$ gates and obeys
\[
 P_{x,y}(c)(z)=
 \begin{cases}y,&z=x,\\c(z),&z\ne x.\end{cases}
\]
The number $r_x$ is independent of the source circuit, and the map is
injective on circuit descriptions.
\end{lemma}
\begin{proof}
Starting from the constant-one wire, form the equality indicator
$e_x(z)=\bigwedge_{j=1}^d\1\{z_j=x_j\}$. A negative literal requires one
NAND gate; conjoining the next literal to the accumulator requires two.
This costs $2d+\#\{j:x_j=0\}\le3d$ gates. If $y=0$, compute
$c(z)\wedge\neg e_x(z)$ in three further NAND gates. If $y=1$, compute
$c(z)\vee e_x(z)$ in three further gates. The original output wire is
available in both cases, including the empty circuit's constant-zero
output. Deleting the final $r_x=2d+\#\{j:x_j=0\}+3$ gates recovers
the original ordered list, proving injectivity. No simplification of dead
or redundant gates is performed.
\end{proof}

For a noiseless history, overwriting the current input with its true label
preserves all earlier labels. This remains true for repeated inputs.
Injectivity is a statement about descriptions, not functions: two
different functions can coincide after their value at one point is
overwritten. Distinguishing these facts matters for prior-mass estimates.

\section{A linear converse for minimum-consistent-circuit prediction}

\begin{theorem}[Deterministic Occam converse]\label{thm:occam}
At every round let the learner predict with any minimum-size circuit
consistent with the preceding observations, with arbitrary tie-breaking.
If it has made $B_t$ mistakes through round $t$, then
\[
 C_{A_t}(f)\le R_d B_t.
\]
In particular a complete reveal of $A$ entails at least
$C_A(f)/R_d$ mistakes, and a complete reveal of the cube entails at least
$C(f)/R_d$ mistakes.
\end{theorem}
\begin{proof}
Initially the minimum consistent size is zero. Adding a constraint cannot
decrease that minimum. If the current prediction is correct, the selected
minimum-size circuit remains consistent, so the minimum does not change.
If the prediction is wrong, Lemma~\ref{lem:repair} produces a consistent
circuit with at most $R_d$ additional gates. Induction proves the result.
\end{proof}

This is a converse, not a polynomial upper bound on the mistakes of
arbitrary minimum-size tie-breaking. The proof does not assert that a
wrong prediction eliminates a substantial fraction of competing circuits.
Exact minimum-circuit selection can be computationally intractable; no
efficient implementation is claimed.

\section{Bayesian sampling, log loss, and complete reveals}

Let $\pi$ be a prior on a countable set of circuit descriptions, with every
finite noiseless history assigned positive mass. Let
\[
 E_t=\{c:c(x_i)=f(x_i)\ (1\le i\le t)\},\qquad
 Z_t=\pi(E_t),\qquad L_t=-\log Z_t.
\]
The posterior is $\pi(\cdot\mid E_{t-1})$. Define
\[
 p_t=\pi\{c(x_t)=f(x_t)\mid E_{t-1}\},\quad q_t=1-p_t,
 \quad m_t=\sum_{i=1}^t q_i.
\]
A posterior-sampling predictor draws a fresh circuit from that posterior
and outputs its label. Its conditional mistake probability is exactly
$q_t$. For a fixed input sequence, $m_t$ is its expected number of
mistakes. For an adaptive or random input sequence, $m_t$ is the random
sum of these conditional probabilities, and expected actual mistakes are
$\E m_t$ (or the analogous stopped/infinite sum).

Irrespective of the query order, exact conditioning gives
\begin{equation}\label{eq:logloss}
 L_t=\sum_{i=1}^t-\log p_i,\qquad m_t\le L_t.
\end{equation}
On a complete reveal of $A$, the final value is the deterministic quantity
\[
 L_A=-\log\pi\{c:c|_A=f|_A\}.
\]
The forward inequality $m_t\le L_t$ does not by itself give a converse:
one incorrect prediction can correspond to an arbitrarily unlikely
correct label and hence arbitrarily large log loss. The next sections
prevent this by repairing descriptions at controlled prior cost.

\section{A natural geometric gate-count prior}

Define $\pi_d$ by first drawing $S$ with
$\Prb(S=s)=2^{-(s+1)}$, $s=0,1,\ldots$. Conditional on $S=s$, draw the
ordered inputs to gate $i$ independently and uniformly from its
$d+i+1$ earlier wires. Thus each particular circuit $c$ of size $s$ has
\begin{equation}\label{eq:prior}
 \pi_d(c)=2^{-(s+1)}\prod_{i=1}^s(d+i+1)^{-2}.
\end{equation}
This is a normalized probability distribution: the product counts the
inverse of the number of wiring lists of length $s$. It is a prior on
descriptions and induces a prior on functions by summing their masses.

For $u\ge0$, put
\begin{equation}\label{eq:A}
 A_d(u)=\log(4/3)+R_d\log2+
          2R_d\log\bigl(u/\log2+4d+6\bigr).
\end{equation}

\begin{lemma}[A lower bound on predictive label mass]\label{lem:predictive}
Under \eqref{eq:prior}, any noiseless history and any next input satisfy
\[
 p_t\ge \exp[-A_d(L_{t-1})].
\]
Consequently every finite history satisfies
\begin{equation}\label{eq:pathwise}
 L_t\le2A_d(L_t)m_t.
\end{equation}
\end{lemma}
\begin{proof}
Set $K=\lceil L_{t-1}/\log2\rceil+1$. Since
$\pi_d(S>K)=2^{-(K+1)}$, the posterior mass on $S>K$ is at most
$2^{-(K+1)}e^{L_{t-1}}\le1/4$. At least $3/4$ of posterior mass
therefore lies on consistent source circuits of size at most $K$.

Repair every such source using $P_{x_t,f(x_t)}$. The resulting descriptions
are distinct, fit all the observations including the new one, and have
prior-to-source probability ratio
\[
 2^{-r_{x_t}}\prod_{i=s+1}^{s+r_{x_t}}(d+i+1)^{-2}
 \ge\bigl[2(d+K+R_d+1)^2\bigr]^{-R_d}.
\]
Summing over the injective image and dividing by the preceding normalizer
gives the lower bound on $p_t$, because
$d+K+R_d+1\le L_{t-1}/\log2+4d+6$.

For $0<b\le p\le1$,
\[
 -\log p\le\frac{-\log b}{1-b}(1-p).
\]
Indeed, $(-\log p)/(1-p)$ decreases with $p$, with the continuous value
one at $p=1$. The function $A_d$ is increasing and exceeds $\log2$.
Use the common lower bound $b=e^{-A_d(L_t)}\le1/2$ for all rounds up to
$t$ and sum. This gives \eqref{eq:pathwise}.
\end{proof}

\begin{theorem}[Geometric circuit-prior converse]\label{thm:bayes}
Suppose posterior sampling with prior \eqref{eq:prior} completely reveals
$A\subseteq\cX$. Let $M$ be its expected total number of mistakes,
including any randomness in a predictable query order. Put $C=C_A(f)$.
Then
\begin{align}
 M&\ge\frac{C\log2}{2A_d(C\log2)},\label{eq:lower}\\
 C&\le\frac{12R_dM}{\log2}
       \log\left(4d+6+\frac{12R_dM}{\log2}\right).
       \label{eq:inverted}
\end{align}
In particular $C=O(dM\log(d+dM))$. Conversely,
\begin{equation}\label{eq:forward}
 M\le(C+1)\log2+2C\log(d+C+1).
\end{equation}
All constants are independent of the reveal order and of $|A|$.
\end{theorem}
\begin{proof}
Any circuit agreeing with $f$ on $A$ has at least $C$ gates, so
$\pi_d(E_A)\le\pi_d(S\ge C)=2^{-C}$ and $L_A\ge C\log2$.
The function $u/A_d(u)$ is increasing on $[0,\infty)$: its derivative
has numerator
\[
 A_d(u)-uA_d'(u)
 =\log(4/3)+R_d\log2+
 2R_d\left[\log v-\frac{v-(4d+6)}v\right]>0,
 \quad v=u/\log2+4d+6\ge10.
\]
Equation~\eqref{eq:pathwise} on each complete-reveal path therefore gives
\[
 m\ge \frac{L_A}{2A_d(L_A)}
     \ge \frac{C\log2}{2A_d(C\log2)}.
\]
Taking expectations proves \eqref{eq:lower}.

Since $C+4d+6\ge10$ and $R_d\ge6$,
$A_d(C\log2)\le3R_d\log(C+4d+6)$. Hence
$C\le a\log(b+C)$ with
$a=6R_dM/\log2$ and $b=4d+6$.
For $a>0,b\ge e$, the implication
\[
 x\le a\log(b+x)\quad\Longrightarrow\quad
 x\le2a\log(b+2a)
\]
follows by setting $x_0=2a\log(b+2a)$, observing
$b+x_0\le(b+2a)^2$, and using the increase of
$x-a\log(b+x)$ for $x\ge x_0$. This proves \eqref{eq:inverted};
if $M=0$, \eqref{eq:lower} already forces $C=0$.

For the upper bound, choose a minimum consistent circuit of size $C$.
Its probability in \eqref{eq:prior} is at least
$2^{-(C+1)}(d+C+1)^{-2C}$. Thus
$m\le L_A\le(C+1)\log2+2C\log(d+C+1)$ on every complete path.
Taking expectations gives \eqref{eq:forward}.
\end{proof}

\begin{remark}[Partial histories]
Before full revelation, \eqref{eq:pathwise} still holds, and applying
the same reasoning gives a small circuit consistent with the observed
points in terms of $m_t$. It bounds $C_{A_t}(f)$, not the unrestricted
$C(f)$. No inference about unobserved labels is implicit.
\end{remark}

\section{A prefix circuit prior with a sharper converse}

The logarithmic factor in Theorem~\ref{thm:bayes} comes from the wire
addresses of appended gates: an exception appended to a longer circuit
costs more prior probability. An explicit exception constructor removes
this dependence. The following is a precise prior, not an assertion that
all circuit encodings have this property.

Use the following prefix language, with $d$ known to the decoder.
\begin{itemize}
\item The string \texttt{00} denotes constant zero.
\item A string \texttt{01}$b$ denotes a NAND circuit with $s\ge1$ gates.
The prefix $b$ starts with the unary length $1^s0$, followed, for each
gate $i$, by two wire addresses of length
$\lceil\log_2(d+i+1)\rceil$. Out-of-range addresses are invalid.
\item A string \texttt{1}$xyp$ denotes the overwrite of program $p$ at
input $x$ with label $y$, where $x$ is encoded in exactly $d$ bits, $y$
in one bit, and $p$ is another valid program.
\end{itemize}
The valid finite strings are prefix-free: the initial tags separate the
cases, unary lengths determine base-code lengths, and the recursive case
has fixed-size metadata before its single prefix-free child. Every
program denotes a Boolean function and every Boolean circuit has such a
description. A program of bit length $\ell$ compiles into a NAND circuit
of at most $3\ell$ gates, by Lemma~\ref{lem:repair} and
$R_d\le3(d+2)$.

For valid strings $p$, define
\begin{equation}\label{eq:prefixprior}
 \widetilde\pi_d(p)=Z_d^{-1}2^{-2\ell(p)},\qquad
 Z_d=\sum_{p\text{ valid}}2^{-2\ell(p)}.
\end{equation}
Kraft's inequality and the code \texttt{00} imply
$1/16\le Z_d\le1$. Thus this is a proper circuit-description prior.
Write $K_A(f)$ for the shortest program length agreeing with $f$ on $A$.

\begin{theorem}[Linear converse with exception codes]\label{thm:prefix}
For posterior sampling with \eqref{eq:prefixprior} and a complete reveal
of $A$, expected mistakes $M$ satisfy
\[
 C_A(f)\le
 \frac{6(d+2)}{1-2^{-2(d+2)}}M+12.
\]
Moreover $M\le2K_A(f)\log2$, and
$K_A(f)=O(C_A(f)\log(d+C_A(f))+1)$. Consequently this same prior has
both an $O(dM+1)$ converse and the usual circuit-size upper bound.
\end{theorem}
\begin{proof}
The map $p\mapsto\texttt{1}xyp$ is injective, preserves earlier consistent
labels when $y=f(x)$, and adds exactly $d+2$ bits. Its prior ratio is
the constant $\beta=2^{-2(d+2)}$. Applying it to all descriptions in
the preceding version space gives predictive correct-label probability
at least $\beta$. Therefore
\[
 L_A\le\frac{2(d+2)\log2}{1-\beta}\,m
\]
on every complete path.

For any real $r\ge0$, Kraft's inequality gives
\[
 \widetilde\pi_d\{\ell(p)\ge r\}
 \le Z_d^{-1}2^{-r}\sum_p2^{-\ell(p)}\le16\,2^{-r}.
\]
Every program agreeing with $f$ on $A$ has length at least $C_A(f)/3$,
since its compiled circuit has at most three times its length in gates.
It follows that
$L_A\ge(C_A(f)/3)\log2-\log16$. Combine the two inequalities and
take expectations for the asserted converse. A shortest fitting program
has probability $2^{-2K_A(f)}/Z_d$, so
$m\le L_A\le2K_A(f)\log2+\log Z_d\le2K_A(f)\log2$.
Finally a circuit of size $s\ge1$ has a base description of length
\[
 s+3+2\sum_{i=1}^s\lceil\log_2(d+i+1)\rceil
 \le s+3+2s\lceil\log_2(d+s+1)\rceil.
\]
For $s=0$, use \texttt{00}. This proves the forward complexity estimate.
\end{proof}

The squared code weight in \eqref{eq:prefixprior} is useful: prefix
freeness alone does not give an exponentially decreasing tail for the
ordinary weights $2^{-\ell(p)}$. The additional length penalty makes
the tail estimate uniform over this language. The result is about the
specified prior, not an invariant property of all ``universal'' priors.

\section{Random input orders and iid sampling}

All results above hold for any fixed reveal order. The Bayesian lower
bounds were established pathwise for the sum of conditional mistake
probabilities, using a final log loss depending only on the revealed
set. Consequently the same bounds hold after averaging over random
orders, including orders selected using previous sampled predictions.
The input for the current round must still precede that round's sample.

\begin{corollary}[Infinite iid input sequence]\label{cor:iid}
Let $D$ be any distribution on $\cX$ and
$A=\{x:D(x)>0\}$. Run any of the specified learners on an infinite iid
input sequence from $D$. Every point of $A$ is eventually seen almost
surely. Theorem~\ref{thm:occam} holds for its total actual mistakes, and
Theorems~\ref{thm:bayes} and \ref{thm:prefix} hold for expected total
posterior-sampling mistakes, with target complexity $C_A(f)$.
When $D$ has full support this is $C(f)$.
\end{corollary}
\begin{proof}
The set $A$ is finite, and each of its points has a finite first-occurrence
time almost surely. After all have appeared, every consistent predictor
is correct on all future draws. Predictions at repetitions are already
correct before this stopping time as well. Thus there are at most $|A|$
possible mistake rounds. Conditional expectation and Tonelli's theorem
identify expected actual mistakes with the expected sum of conditional
mistake probabilities; the complete-reveal results then apply.
\end{proof}

No lower bound on the positive probabilities $D(x)$ is required for
total mistakes. Such a bound would matter for how long it takes to
reveal the support: very rare inputs can postpone the decisive mistakes
for arbitrarily many rounds. A small asymptotic error rate or a small
number of mistakes over a finite horizon is therefore a different claim
from a small total number of mistakes over the infinite sequence.

\section{Limits of the converse}

\paragraph{Arbitrary learners.}
For any prescribed function $f$, a learner that always returns $f(x)$
makes zero mistakes. By counting NAND circuits, there are Boolean
functions with circuit size of order $2^d/d$. Thus no bound in terms of
mistakes and $d$ alone holds for arbitrary learners without restricting
their initial knowledge or representation. The same observation applies
to an arbitrarily selected prior: a point mass on a complex $f$ makes
zero posterior-sampling mistakes. Even full support does not repair this.
For example, assigning mass $1-\eta$ to $f$ and spreading the remaining
$\eta$ over all Boolean functions gives cumulative sampling mistakes at
most $-\log(1-\eta)$, however complex $f$ is.

\paragraph{Unrevealed inputs.}
For $d\ge2$, a learner can observe only a point where $f=0$, while the
remaining truth table has almost arbitrary circuit complexity. Taking
the input distribution to be concentrated on that point gives no control
of global $C(f)$. The valid complexity in the iid theorem is the minimum
size needed to match $f$ on the distribution's support.

\paragraph{Posterior majority.}
Sampling mistake probability is $1-p_t$, whereas majority prediction is
correct whenever $p_t>1/2$. A sequence with $p_t=3/4$ on every round has
no majority mistakes but positive sampling mistakes and log loss on every
round. This can occur for a prior with independently biased truth-table
bits. The mass-repair proofs therefore do not establish a majority
converse, even though they establish one for posterior sampling. The
separate deterministic Occam theorem concerns minimum-size consistent
circuits, not posterior majorities.

\paragraph{A minimum description versus a function prior.}
The geometric and prefix priors here assign mass to descriptions, and
multiple descriptions of the same function are all retained. Injectivity
of repairs makes their masses add correctly. For a prior specified
directly on functions, the overwrite map has up to two preimages and
its weights must be checked separately. Merely declaring that the prior
``prefers smaller circuits'' does not imply the quantitative ratios or
tail bounds used above.

\paragraph{Depth, noise, and computation.}
The conclusions count all NAND gates and allow their depth to grow.
One-point overwrites do not preserve a fixed depth architecture. Repeated
conflicting labels are excluded; without realizability the consistent
posterior can be empty. These are statistical and representation
statements, with no efficient circuit minimization or posterior sampler
asserted.

\paragraph{Acknowledgment and status.}
The argument develops the one-point exception-repair and log-loss ideas
in the collaborative circuit-induction discussion. The explicit circuit
priors, injectivity checks, and converse estimates are supplied here.
This manuscript has not been formalized in Lean.
\end{document}
